\documentclass[10pt,a4paper]{amsart}
\usepackage[margin=19mm]{geometry}
\usepackage{amsmath,amssymb,mathtools}
\usepackage[scr=rsfs,cal=euler]{mathalfa}
\usepackage{xfrac}
\usepackage[unicode,hidelinks,bookmarks=false]{hyperref}
\newcommand{\ie}{i.e.\ }
\newcommand{\cf}{cf.\ }
\newcommand{\resp}{respectively\ }
\newcommand{\Subsection}{\subsection}
\providecommand{\nobreakdash}{\nobreak\hbox{-}}
\setlength{\emergencystretch}{2em}
\multlinegap0pt
\usepackage{amssymb}
\usepackage{mathtools}
\usepackage{csquotes}
\usepackage[shortlabels]{enumitem}
\usepackage{tikz}
\usepackage{float}
\usepackage{xcolor}
\usepackage{xfrac}
\newtheorem{proposition}{Proposition}
\newtheorem{lemma}{Lemma}
\newcommand{\C}{\mathcal C}
\def \GL{\sf{GL}}
\newcommand{\K}{\mathbb K}
\newcommand{\R}{\mathbb R}
\title{{\small\textbf{A CRITERION FOR TITS ALTERNATIVE ON THE CENTRALIZER OF A MATRIX}}}
\author{Adem Zeghib}
\date{Source: arXiv:2604.06968v1 (CC BY 4.0)}
\begin{document}
\maketitle

\noindent\emph{Source and licence.} {\small\textbf{A CRITERION FOR TITS ALTERNATIVE ON THE CENTRALIZER OF A MATRIX}}. Version arXiv:2604.06968v1 (CC BY 4.0), distributed by arXiv under the Creative Commons Attribution 4.0 International licence, http://creativecommons.org/licenses/by/4.0/, as recorded in the arXiv licence metadata for this exact version and shown on the abstract page https://arxiv.org/abs/2604.06968v1. Full licence notice: \texttt{licenses/arxiv-2604.06968-CC-BY-4.0.txt}.\par\smallskip

\noindent\textit{Editorial scope.} Counted unit: the article's proposition prop4 (source0.tex lines 221--223). The article's complete original proof of it is delivered with it (source0.tex lines 225--303, one \texttt{proof} environment of 79 lines). No mathematics is written, repaired or re-derived: the statement and the proof are the article's own lines, in the article's own notation, and no retained line is substituted or re-flowed.  No further source-local context is needed: every cross-reference of every retained line resolves inside the two delivered spans.  Nothing was omitted from any retained span, so the package carries no omission record.  No retained line contains a citation, so the wrapper carries no bibliography and no bibliography entry.  No retained line contains a figure environment, an image-inclusion command or drawing code, so no image is delivered and the package declares no asset dependency.  Editorial formatting only, in addition: an AMS \texttt{amsart} preamble that restates the article's own declarations and macros used by the retained lines, namely the environments \texttt{proposition}, \texttt{lemma} on separate counters, together with the standard packages those lines need.  The retained environments print in this document as Proposition 1, Lemma 1 in order of appearance, whereas the article prints the same items with the numbering of its own full document.  The complete native source of the article as distributed in the arXiv e-print for this exact version is bundled unchanged as \texttt{sources/198057/source0.tex}.
\begin{proposition}\label{prop4}
If $T \in \GL(n,\K)$ has two Jordan blocks of the same size for the same eigenvalue, then $\C_{\GL(n,\overline{\K})}(T)$ contains a copy of $\GL(2,\overline{\K})$. Moreover, if $\K=\R$, $\C_{\GL(n,\R)}(T)$ contains a copy of $\GL(2, \R)$.
\end{proposition}

\begin{proof}
The matrix $T$ is conjugate in $\GL(n,\overline{K})$ to the following matrix
    $T \sim_{\GL(n,\overline{\K})}
J'=\begin{pmatrix}
J & 0& 0 \\
0 & J & 0\\
0 & 0 & A 
\end{pmatrix} $
then 
$\begin{pmatrix}
aI & bI & 0 \\
cI & dI & 0\\
0 & 0 & I 
\end{pmatrix} $
commutes with $J'$ for every $a,b,c,d \in \overline{\K} $ and is invertible as long as $ad-bc\ne 0$.
The following map is an injective homomorphism : 

   \[ \Psi  :
\begin{array}{ccccc}
&  \GL(2,\overline{\K})  & \longrightarrow & \GL(n,\overline{\K}) & \\
& \begin{pmatrix}
a & b  \\
c & d 
\end{pmatrix}  & \longmapsto &
\begin{pmatrix}
aI & bI & 0 \\
cI & dI & 0 \\
0 & 0 & I
\end{pmatrix} \\
\end{array}
 \]



If $\K=\R$, we need the following lemma. 

\begin{lemma}\label{lemma5}
    Let $T$ be a real matrix such that $T$ is $\mathbb{C}$-conjugate to $diag(J(\lambda),A )$, with $J(\lambda)$ a Jordan block of size $s$. \newline
    If $\lambda \notin \R$ then
    \begin{itemize}
        \item $T$ is $\mathbb{C}$-conjugate to $diag(J(\lambda), J(\overline{\lambda})), A')$, with $J(\overline{\lambda})$ also of  size $s$.
        \item $T$ is $\R$-conjugate to $diag(N,B)$ with $N$ and $B$ real matrices.
    \end{itemize}
\begin{proof}
    First, since $T$ is a real matrix, then $\overline{T}=T$, which leads to the existence of a Jordan block $J(\overline{\lambda})$ of the same size.
    Then, we can easily check that 
$$\begin{pmatrix}
I_s & I_s \\
iI_s & -iI_s  
\end{pmatrix}^{-1}= \frac{1}{2}\begin{pmatrix}
I_s & -iI_s \\
I_s & iI_s  
\end{pmatrix}$$ 

and that the conjugate

$$\begin{pmatrix}
I_s & I_s \\
iI_s & -iI_s  
\end{pmatrix}
\begin{pmatrix}
J(\lambda) & 0 \\
0 & J(\overline{\lambda})  
\end{pmatrix}
\begin{pmatrix}
I_s & -iI_s \\
I_s & iI_s  
\end{pmatrix} \quad \text{is a real matrix}$$

We can continue in the same fashion and gather Jordan blocks of the same size that are conjugate to obtain the matrix $B$. Therefore, $T$ and $diag(N,B)$ are $\mathbb{C}$-conjugate. We conclude using the classical fact that if two real matrices are $\mathbb{C}$-conjugate they are also $\R$-conjugate.

\end{proof}
    
\end{lemma}


We continue our proof. Let $T$ be a real matrix as in the statement of Proposition \ref{prop4}.  Using the last lemma we get that $T$ is $\R$-conjugate to $diag(N,N,A')$ with $N$ and $A$ real matrices. We conclude using the homomorphism $\Psi$ as before to show that there is a copy of $\GL(2,\R)$ in $\C_{\GL(n,\R)}(T)$.

\end{proof}                                   

\end{document}
